% Source: https://artofproblemsolving.com/community/c3h3352781_interesting_counting_problem
% Stable identity: aops-topic-3352781-q01
The number of $5$s must be even. Four $5$s give one outcome. Exactly two give $\binom42=6$ position choices; the remaining ordered pair from $\{1,2,3,4,6\}$ has square product in $7$ ways (four pairs from $\{1,4\}$, and $(2,2),(3,3),(6,6)$), giving $42$. With no $5$s, let the square class be $\{1,4\}$ and the three nonsquare classes be $\{2\},\{3\},\{6\}$. The valid class patterns are: four square entries ($16$); two square entries and two from one nonsquare class ($3\cdot6\cdot4=72$); four equal nonsquare entries ($3$); two each from two nonsquare classes ($3\cdot6=18$); and one from each nonsquare class with one square entry ($24\cdot2=48$). These exhaust the even prime-exponent possibilities, totaling $157$. Therefore there are $1+42+157=200$ favorable outcomes, probability $\dfrac{200}{1296}=\dfrac{25}{162}$, and $\boxed{m+n=187}$.
